% Copyright (c) 2026 Geovana Neves. Licensed under CC BY 4.0 (https://creativecommons.org/licenses/by/4.0/). See LICENSE-AND-AUTHORSHIP.md.
%
\section{\code{per\_blade}}

\begin{frame}{One shared window, one row per blade}
\begin{itemize}\small
  \item \code{probes/<point>\_per\_blade\_<ALIAS>.csv} averages each blade
        over the last part of the run: the matrix row's averaging window.
  \item Every blade shares that window. Each row states that blade's start
        and end azimuth.
  \item The window is the one the unsteady plots and POLAR use, from
        \code{LAST\_REVS\_AVG} or \code{LAST\_ITERS\_AVG} on the matrix row.
  \item There is one row per blade.
\end{itemize}
\end{frame}

\begin{frame}{The same time interval makes blades comparable}
\begin{itemize}
  \item Each blade starts at a different azimuth. Its azimuth columns
        state that difference explicitly.
  \item A separate passage window for each blade would also sample each
        blade at a different time, mixing azimuth and history differences.
  \item Averaging from step one mixes the transient with the answer.
\end{itemize}
\begin{takeaway}
The last shared window keeps the sampling time the same for every blade.
\end{takeaway}
\end{frame}

\begin{frame}{The leading columns of each blade row}
\small Each row leads with these fields, in this order:
\begin{enumerate}
  \item \code{REDUCTION}, \code{ROTOR}, \code{BLADE}, \code{FAMILY}.
  \item The window: \code{FIRST\_STEP}, \code{LAST\_STEP}, \code{STEPS}.
  \item \code{AZIMUTH\_START}, \code{AZIMUTH\_END}.
  \item The condition block, then the moment point.
\end{enumerate}
\end{frame}

\begin{frame}{Blade plot columns share one heading}
\begin{itemize}\small
  \item A blade's columns are plots named for its family.
        A plot is \code{<parameter>\_<group>}; a group cut per blade is
        named for the blade's family.
  \item Blade one's \code{CL\_MRP\_Blade1} and \code{FX\_LOCAL\_Blade1}
        become \code{CL\_MRP} and \code{FX\_LOCAL} in its row.
        Removing the family makes the blades line up under one heading.
  \item The pproc requests these with \code{families = "each"} or
        \code{frame = "LOCAL\_AXIS"} over the rotor.
\end{itemize}
\end{frame}

\begin{frame}{Start and end azimuth belong to that blade}
\small At the first and last step of the shared window, use
\[
\begin{split}
\mathrm{azimuth}=\biggl(&\mathrm{blade1.azimuth\_deg}
 +\mathrm{blade\ position}\,\frac{360}{\mathrm{blades}}\\
 &+\mathrm{sense}\,\mathrm{step}\,
 \frac{360}{\mathrm{steps\_per\_revolution}}\biggr)\bmod 360.
\end{split}
\]
\begin{itemize}
  \item \code{blades} is the rotor's count. A periodic sector carrying
        two blades of four spaces them a quarter turn apart.
  \item Without the rotor's clock, the azimuths are \code{NA}.
        The averages are still written.
\end{itemize}
\end{frame}

\begin{frame}{The reference identifies the blade families}
\begin{itemize}
  \item The rotor block's \code{families\_blades} in the reference artifact
        identifies its blades.
  \item The run also records these families.
  \item A row using flat rotor keys and citing no block gets no per-blade
        table. \code{products.json} states why.
\end{itemize}
\end{frame}
